\documentclass[11pt]{article}
\usepackage{icmlko}
\begin{document}

\title{수집 가격표 --- 축 하나가 검색 시계열 전체보다 비싸고, 우리가 세 노트를 쓴 축은 공짜다}
\author{월드모델 실험실}
\date{노트 496 · 2026년 8월 3일}
\maketitle

\begin{abstract}
노트 490 이 시계열 수집에 판 $-0.0283$ 을 매겼다. 같은 자를 \textbf{속성
수집 일곱}에도 대서 \textbf{전 수집 가격표}를 만들었다(짝 씨앗 $12$ ·
학습과 예보 양쪽에서 통째로 없앰).

\begin{center}
\begin{tabular}{lrrl}
\hline
수집 & 판 차 & 부호 & 갈래 \\
\hline
\texttt{target\_breadth} & $\mathbf{-0.0306}$ & $12/12$ & 속성 \\
시계열 둘 다(트렌드$+$위키) & $-0.0283$ & $12/12$ & 시계열 \\
태그 \texttt{tag\_c*} & $-0.0236$ & $12/12$ & 속성 \\
트렌드 셋 & $-0.0206$ & $12/12$ & 시계열 \\
\texttt{venue\_prominence} & $-0.0134$ & $12/12$ & 속성 \\
위키 셋 & $-0.0071$ & $12/12$ & 시계열 \\
\texttt{goods\_scale} & $-0.0065$ & $12/12$ & 속성 \\
\texttt{media\_push} & $-0.0030$ & $11/12$ & 속성 \\
\texttt{gen} & $-0.0025$ & $11/12$ & 속성 \\
\texttt{grp} & $-0.0011$ & $9/12$ & 속성 \\
\textbf{\texttt{entry\_friction}} & $\mathbf{-0.0009}$ & $8/12$ & 속성 \\
\hline
\end{tabular}
\end{center}

두 가지가 크게 걸린다.
\end{abstract}

\begin{figure}[t]
\centering
\includegraphics[width=\textwidth]{figs/f1.pdf}
\caption{\textbf{가} --- 가격표. \textbf{나} --- 두 자 견주기.}
\end{figure}

\section{하나 --- 축 하나가 수집 전체보다 비싸다}

\texttt{target\_breadth} 를 빼면 판이 $-0.0306$ 내려간다. \textbf{검색
시계열 여섯 축 전부를 빼는 것($-0.0283$)보다 크다.} 날짜 통제 판에서도
$-0.0377$ 로 제일 크다.

이 축은 사람이 붙이는 속성 표지 하나다. 두 해 동안 이 실험실은
\textbf{그것이 얼마짜리인지 한 번도 안 쟀다.} 모형을 $0.004$ 씩 밀면서.

\section{둘 --- 세 노트를 쓴 축이 공짜다}

\texttt{entry\_friction} 은 $-0.0009$ 이고 부호가 $8/12$ 다. \textbf{잡음
안이다.}

이 축은 노트 $441 \sim 443$ 의 주인공이었다 --- 공유 축 다섯 중 유일하게
도메인 간 부호가 안 맞는다고, 집 밖에서 가리면 KR 이 $+0.2330 \to
+0.5504$ 로 뛴다고, 라벨 다섯이면 뒤집어 쓸 수 있다고. 노트 463 이 그
셋을 축소 모형 산물로 되돌렸고, 노트 $465 \sim 468$ 이 다시 넷을 더 썼다.
\textbf{일곱 노트다.}

\begin{quote}
그 축은 판에서 $0.0009$ 짜리다.
\end{quote}

집 밖 전이에서 크게 보였던 것은 지금도 사실이다(노트 $465$: 두 수집이
진짜로 반대를 원한다). 다만 \textbf{집안에서 그 축이 얼마인지를 먼저
알았다면 일곱 노트를 그렇게 안 썼을 것이다.}

\begin{quote}
\textbf{조항} --- \textbf{어떤 축을 두 노트 넘게 파기 전에 그 축의 판
값어치를 먼저 잰다.} 값어치가 잡음 안이면 그 축에 대한 발견은 크기가
아니라 \emph{방향}만 말할 수 있다.
\end{quote}

\section{사전 등록은 못 가렸다}

``격차 몫(노트 489)과 판 값어치의 순위가 다르다''에 걸었다. 넷을 짝지어
보니 스피어만 $0.800$ 으로 \emph{같은 순위}처럼 보이는데, \textbf{짝이
넷이라 $p \approx 0.33$ 이다.} 통과도 반증도 아니고 \textbf{못 가른다.}
노트 495 조항을 스스로에게 적용한다 --- 자의 잡음보다 작은 것을 재려
했다.

\section*{남은 것}

\begin{itemize}
\item \textbf{\texttt{target\_breadth} 는 무엇인가} --- 제일 비싼
수집인데 이 실험실은 그 안을 들여다본 적이 없다. 값 분포 · 덮음 ·
누가 붙였나.
\item \textbf{태그 $-0.0236$} --- 둘째로 비싼 속성 수집. 도메인별
\texttt{tag\_c*} 는 서로 다른 것을 세는데 한 덩어리로 쟀다.
\item \textbf{가격표는 \emph{빼는} 값이다} --- \emph{더하는} 값이 아니다.
덮음이 낮은 축은 싸 보일 수밖에 없다(노트 493 의 단서).
\end{itemize}

\end{document}
